%%
%% part2-book-components/ch10-code.tex
%%
%% Chapter 10: Code and Algorithms — Syntax-highlighted
%% code with minted, and pseudocode with algorithm.
%%

\chapter{کد و الگوریتم}
\label{chap:code}

\section{چرا این موضوع مهم است؟}
\label{sec:code-why}

در کتاب‌های برنامه‌نویسی، نمایش حرفه‌ای کد
اهمیت ویژه‌ای دارد. کد باید با رنگ‌آمیزی نحوی
مناسب، فونت خوانا، شماره‌گذاری منظم و کپشن‌های
گویا نمایش داده شود.

\lr{MatinBook} از بسته‌ی \lr{minted} برای نمایش
کد و از بسته‌ی استاندارد \lr{algorithm} برای
شبه‌کد استفاده می‌کند. در این فصل، با هر دو
آشنا می‌شوید و یاد می‌گیرید که چگونه کدهای
فارسی و لاتین را با هم ترکیب کنید.

\mnote{کدها با \lr{minted} و \lr{Pygments} رنگ‌آمیزی می‌شوند که از بیش از ۳۰۰ زبان برنامه‌نویسی پشتیبانی می‌کند.}

\section{نمایش کد با \lr{minted}}
\label{sec:code-minted}

\subsection{بلوک کد پایه}
\label{sec:code-basic}

برای نمایش یک بلوک کد، از محیط \env{minted}
داخل محیط \env{latin} استفاده کنید:

\begin{latin}
\begin{minted}{latex}
\begin{latin}
\begin{minted}{python}
def factorial(n: int) -> int:
    """Calculate n! recursively."""
    if n <= 1:
        return 1
    return n * factorial(n - 1)
\end{minted}
\end{latin}
\codecaption{|\pc{تابع فاکتوریل بازگشتی}|}
\label{code:factorial}
\end{minted}
\end{latin}

خروجی:

\begin{latin}
\begin{minted}{python}
def factorial(n: int) -> int:
    """Calculate n! recursively."""
    if n <= 1:
        return 1
    return n * factorial(n - 1)
\end{minted}
\end{latin}
\codecaption{تابع فاکتوریل بازگشتی}
\label{code:factorial-ch10}

\mnote{محیط \lr{minted} باید همیشه داخل \lr{\textbackslash begin\{latin\}} قرار گیرد. این کار جهت کد را لاتین نگه می‌دارد.}

\subsection{کپشن کد}
\label{sec:code-caption}

برای افزودن کپشن به کد، از دستور \cmd{codecaption}
استفاده کنید:

\begin{latin}
\begin{minted}{latex}
\codecaption{|\pc{کپشن کد}|}
\label{code:my-code}
\end{minted}
\end{latin}

\mnote{دستور \lr{\textbackslash codecaption} یک شمارنده‌ی مستقل برای کدها ایجاد می‌کند که به‌صورت «فصل.شماره» نمایش داده می‌شود.}

\subsection{ارجاع به کد}
\label{sec:code-reference}

برای ارجاع به کد، از دستور \cmd{cref} استفاده
کنید:

\begin{latin}
\begin{minted}{latex}
|\pc{همان‌طور که در}|\cref{code:factorial}|\pc{می‌بینیم...}|%
\end{minted}
\end{latin}

خروجی: همان‌طور که در \cref{code:factorial-ch10}
می‌بینیم...

\mnote{دستور \lr{\textbackslash cref} به‌طور خودکار نوع مرجع را تشخیص می‌دهد و پیشوند «کد» را اضافه می‌کند.}

\section{کامنت فارسی در کد}
\label{sec:code-persian-comments}

یکی از ویژگی‌های مهم \lr{MatinBook}، پشتیبانی
از کامنت فارسی در کد است. برای این کار، از دو
دستور استفاده می‌کنیم:

\begin{itemize}
    \item \cmd{pcc\{...\}}: برای کامنت فارسی
    در پایتون (و زبان‌های مشابه که \lr{\#}
    دارند).
    \item \cmd{pc\{...\}}: برای متن فارسی غیرکامنت
    (مثل \lr{LaTeX}، \lr{bash}).
\end{itemize}

\mnote{دستور \lr{\textbackslash pcc} علامت \lr{\#} را به‌طور خودکار اضافه می‌کند. دستور \lr{\textbackslash pc} فقط متن فارسی را با فونت مناسب نمایش می‌دهد.}

\subsection{کامنت فارسی در پایتون}
\label{sec:code-pcc}

برای کامنت فارسی در پایتون، از \cmd{pcc} داخل
\lr{|...|} استفاده کنید:

\begin{latin}
\begin{minted}{python}
x = 5  |\pcc{مقدار متغیر}|
y = 10  |\pcc{مقدار دیگر}|
result = x + y  |\pcc{جمع دو متغیر}|
\end{minted}
\end{latin}


خروجی:

\begin{latin}
\begin{minted}{python}
x = 5  |\pcc{مقدار متغیر}|
y = 10  |\pcc{مقدار دیگر}|
result = x + y  |\pcc{جمع دو متغیر}|
\end{minted}
\end{latin}

\mnote{اگر به‌جای \lr{\textbackslash pcc} از \lr{\# |\textbackslash pc\{...\}|} استفاده کنید، \lr{Pygments} علامت \lr{\#} را به‌عنوان شروع کامنت می‌بیند و \lr{|} را نادیده می‌گیرد. نتیجه، مربع‌های خالی است.}

\subsection{متن فارسی در \lr{LaTeX}}
\label{sec:code-pc}

برای متن فارسی غیرکامنت (مثلاً در \lr{LaTeX})،
از \cmd{pc} استفاده کنید:

\begin{latin}
\begin{minted}{latex}
\begin{latin}
\begin{minted}{latex}
\documentclass{matinbook}
\begin{document}
\chapter{|\pc{فصل نمونه}|}
\end{document}
\end{minted}
\end{latin}
\end{minted}
\end{latin}

خروجی:

\begin{latin}
\begin{minted}{latex}
\documentclass{matinbook}
\begin{document}
\chapter{|\pc{فصل نمونه}|}
\end{document}
\end{minted}
\end{latin}

\mnote{دستور \lr{\textbackslash pc} باید همیشه داخل \lr{|...|} قرار گیرد. اگر بیرون از \lr{|...|} باشد، به‌عنوان متن خام نمایش داده می‌شود.}

\subsection{جدول تصمیم‌گیری}
\label{sec:code-decision-table}

جدول \cref{tab:code-persian} به شما کمک می‌کند
که در هر موقعیت، از کدام دستور استفاده کنید.

\begin{matintable}{انتخاب دستور مناسب برای فارسی در کد}{tab:code-persian}
\begin{tblr}{
    width = \textwidth,
    colspec = {X[l] X[c] X[c]},
    hlines, vlines,
    colsep = 6pt,
    rowsep = 4pt,
    row{1} = {font=\bfseries},
}
موقعیت & زبان & دستور \\
کامنت فارسی & پایتون & \lr{|\textbackslash pcc\{...\}|} \\
کامنت فارسی & \lr{C++}، \lr{Java} & \lr{|\textbackslash pcc\{...\}|} \\
کامنت فارسی & \lr{LaTeX}، \lr{bash} & \lr{|\textbackslash pc\{\textbackslash\# ...\}|} \\
متن فارسی غیرکامنت & همه & \lr{|\textbackslash pc\{...\}|} \\
کد کاملاً لاتین & همه & بدون \lr{\textbackslash pc} \\
\end{tblr}
\end{matintable}

\section{کد درون‌خطی}
\label{sec:code-inline}

برای ارجاع به یک تابع یا متغیر در متن، از
دستور \cmd{inlcode} استفاده کنید:

\begin{latin}
\begin{minted}{latex}
|\pc{تابع}|\inlcode{factorial()}|\pc{یک تابع بازگشتی است.}|%
\end{minted}
\end{latin}

خروجی: تابع \inlcode{factorial()} یک تابع
بازگشتی است.

\mnote{دستور \lr{\textbackslash inlcode} فونت مونواسپیس را در متن اعمال می‌کند و برای کدهای کوتاه (کمتر از ۴۰ کاراکتر) مناسب است.}

\section{الگوریتم‌ها}
\label{sec:code-algorithms}

\subsection{محیط \lr{algorithm}}
\label{sec:code-algorithm-basic}

برای نمایش شبه‌کد، از بسته‌ی استاندارد
\lr{algorithm} استفاده می‌کنیم. این محیط
باید داخل \env{latin} قرار گیرد:

\begin{latin}
\begin{minted}{latex}
\begin{latin}
\begin{algorithm}
\caption{|\pc{جستجوی دودویی}|}
\label{alg:binary}
\begin{algorithmic}[1]
    \Require Sorted array $A$ of $n$ elements
    \Require Target value $x$
    \Ensure Index of $x$ in $A$, or $-1$
    \State $left \gets 0$, $right \gets n - 1$
    \While{$left \leq right$}
        \State $mid \gets (left + right) / 2$
        \If{$A[mid] = x$}
            \State \Return $mid$
        \ElsIf{$A[mid] < x$}
            \State $left \gets mid + 1$
        \Else
            \State $right \gets mid - 1$
        \EndIf
    \EndWhile
    \State \Return $-1$
\end{algorithmic}
\end{algorithm}
\end{latin}
\end{minted}
\end{latin}

خروجی:

\begin{latin}
\begin{algorithm}
\caption{جستجوی دودویی}
\label{alg:binary-ch10}
\begin{algorithmic}[1]
    \Require Sorted array $A$ of $n$ elements
    \Require Target value $x$
    \Ensure Index of $x$ in $A$, or $-1$
    \State $left \gets 0$, $right \gets n - 1$
    \While{$left \leq right$}
        \State $mid \gets (left + right) / 2$
        \If{$A[mid] = x$}
            \State \Return $mid$
        \ElsIf{$A[mid] < x$}
            \State $left \gets mid + 1$
        \Else
            \State $right \gets mid - 1$
        \EndIf
    \EndWhile
    \State \Return $-1$
\end{algorithmic}
\end{algorithm}
\end{latin}

\mnote{عنوان الگوریتم (\lr{\textbackslash caption}) می‌تواند فارسی باشد، ولی متن داخل \lr{algorithmic} باید لاتین باشد.}

\subsection{کامنت فارسی در الگوریتم}
\label{sec:code-algorithm-comment}

برای کامنت فارسی در الگوریتم، از \cmd{pc} داخل
\lr{|...|} استفاده کنید:

\begin{latin}
\begin{minted}{latex}
\begin{algorithmic}[1]
    \Require Sorted array $A$
    \While{$left \leq right$}
        \State $mid \gets (left + right) / 2$
        \If{$A[mid] = x$}
            \State \Return $mid$  \Comment{|\pc{یافتن مقدار}|}
        \EndIf
    \EndWhile
\end{algorithmic}
\end{minted}
\end{latin}

خروجی:

\begin{latin}
\begin{algorithm}
\caption{جستجو با کامنت فارسی}
\label{alg:binary-comment}
\begin{algorithmic}[1]
    \Require Sorted array $A$
    \State $left \gets 0$, $right \gets n - 1$
    \While{$left \leq right$}
        \State $mid \gets (left + right) / 2$
        \If{$A[mid] = x$}
            \State \Return $mid$  \Comment{|\pc{یافتن مقدار}|}
        \ElsIf{$A[mid] < x$}
            \State $left \gets mid + 1$  \Comment{|\pc{نیمه‌ی راست}|}
        \EndIf
    \EndWhile
\end{algorithmic}
\end{algorithm}
\end{latin}

\mnote{کامنت‌های الگوریتم با \lr{\textbackslash Comment\{|\textbackslash pc\{...\}|\}} نوشته می‌شوند.}

\subsection{دستورات \lr{algorithmic}}
\label{sec:code-algorithm-commands}

جدول \cref{tab:algorithm-commands} دستورات
اصلی \env{algorithmic} را خلاصه می‌کند.

\begin{matintable}{دستورات اصلی \lr{algorithmic}}{tab:algorithm-commands}
\begin{tblr}{
    width = \textwidth,
    colspec = {l X[l]},
    hlines, vlines,
    colsep = 8pt,
    rowsep = 4pt,
    row{1} = {font=\bfseries},
}
دستور & کاربرد \\
\lr{\textbackslash Require} & پیش‌نیاز الگوریتم \\
\lr{\textbackslash Ensure} & خروجی الگوریتم \\
\lr{\textbackslash State} & یک دستور ساده \\
\lr{\textbackslash Statex} & دستور بدون شماره \\
\lr{\textbackslash If} / \lr{\textbackslash ElsIf} / \lr{\textbackslash Else} & شرط \\
\lr{\textbackslash EndIf} & پایان شرط \\
\lr{\textbackslash For} / \lr{\textbackslash EndFor} & حلقه‌ی \lr{for} \\
\lr{\textbackslash While} / \lr{\textbackslash EndWhile} & حلقه‌ی \lr{while} \\
\lr{\textbackslash Function} / \lr{\textbackslash EndFunction} & تعریف تابع \\
\lr{\textbackslash Call\{name\}\{args\}} & فراخوانی تابع \\
\lr{\textbackslash Return} & بازگشت مقدار \\
\lr{\textbackslash Comment\{...\}} & کامنت \\
\end{tblr}
\end{matintable}

\subsection{ارجاع به الگوریتم}
\label{sec:code-algorithm-reference}

برای ارجاع به الگوریتم، از \cmd{cref} استفاده
کنید:

\begin{latin}
\begin{minted}{latex}
|\pc{الگوریتم}|\cref{alg:binary}|\pc{فرآیند جستجو را نشان می‌دهد.}|%
\end{minted}
\end{latin}

خروجی: الگوریتم \cref{alg:binary-ch10} فرآیند
جستجو را نشان می‌دهد.

\mnote{برای فهرست الگوریتم‌ها، از \lr{\textbackslash listofalgorithms} در پیش‌متن استفاده کنید.}

\section{تنظیمات پیش‌فرض \lr{minted}}
\label{sec:code-settings}

\lr{MatinBook} تنظیمات زیر را برای \lr{minted}
اعمال می‌کند:

\begin{itemize}
    \item \textbf{فونت:} \lr{JetBrains Mono} با
    اندازه‌ی \lr{footnotesize}.
    \item \textbf{شماره خطوط:} فعال
    (\lr{linenos=true}).
    \item \textbf{شکستن خطوط:} فعال
    (\lr{breaklines=true}).
    \item \textbf{کادر:} خطوط بالا و پایین
    (\lr{frame=lines}).
    \item \textbf{پس‌زمینه:} \lr{matinlightgray!30}.
    \item \textbf{اندازه Tab:} ۴ فاصله
    (\lr{tabsize=4}).
    \item \textbf{حذف تورفتگی مشترک:} فعال
    (\lr{autogobble=true}).
    \item \textbf{استایل رنگ‌آمیزی:} \lr{friendly}.
\end{itemize}

\mnote{برای شخصی‌سازی این تنظیمات، فایل \lr{mb-code.sty} را ویرایش کنید.}

\section{زبان‌های پشتیبانی‌شده}
\label{sec:code-languages}

\lr{minted} از بیش از ۳۰۰ زبان برنامه‌نویسی
پشتیبانی می‌کند. چند زبان پرکاربرد:

\begin{itemize}
    \item \textbf{برنامه‌نویسی:} \lr{python}،
    \lr{cpp}، \lr{c}، \lr{java}،
    \lr{javascript}، \lr{typescript}،
    \lr{rust}، \lr{go}.
    \item \textbf{نشانه‌گذاری:} \lr{html}،
    \lr{xml}، \lr{markdown}، \lr{latex}،
    \lr{yaml}، \lr{json}.
    \item \textbf{پوسته:} \lr{bash}، \lr{zsh}،
    \lr{powershell}.
    \item \textbf{پایگاه داده:} \lr{sql}،
    \lr{postgresql}، \lr{mysql}.
    \item \textbf{سایر:} \lr{text} (ساده)،
    \lr{diff}، \lr{makefile}، \lr{docker}.
\end{itemize}

\mnote{برای دیدن فهرست کامل زبان‌ها، از \lr{pygmentize -L lexers} استفاده کنید.}

\section{خطاهای رایج}
\label{sec:code-mistakes}

\subsection{فراموش کردن \env{latin}}
\label{sec:code-mistake-latin}

اگر کد را داخل \env{latin} قرار ندهید، ممکن
است جهت کد به‌هم بریزد و شماره‌های خطوط
جابه‌جا شوند:

\begin{latin}
\begin{minted}{latex}
% |\pc{اشتباه}|:
\begin{minted}{python}
x = 5  |\pcc{مقدار}|
\end{minted}

% |\pc{درست}|:
\begin{latin}
\begin{minted}{python}
x = 5  |\pcc{مقدار}|
\end{minted}
\end{latin}
\end{minted}
\end{latin}

\subsection{فراموش کردن \cmd{pcc} در پایتون}
\label{sec:code-mistake-pcc}

اگر در پایتون از \lr{\# |\textbackslash pc\{...\}|}
استفاده کنید، \lr{Pygments} علامت \lr{\#} را
به‌عنوان شروع کامنت می‌بیند و \lr{|} را نادیده
می‌گیرد:

\begin{latin}
\begin{minted}{latex}
% |\pc{اشتباه}|:
x = 5  # |\pc{مقدار متغیر}|

% |\pc{درست}|:
x = 5  |\pcc{مقدار متغیر}|
\end{minted}
\end{latin}

\subsection{فراموش کردن \cmd{pc} در \lr{LaTeX}}
\label{sec:code-mistake-pc}

اگر در \lr{LaTeX} متن فارسی را بدون \cmd{pc}
بنویسید، نتیجه به‌صورت مربع‌های خالی نمایش
داده می‌شود:

\begin{latin}
\begin{minted}{latex}
% |\pc{اشتباه}|:
\chapter{فصل نمونه}

% |\pc{درست}|:
\chapter{|\pc{فصل نمونه}|}
\end{minted}
\end{latin}

\section{تمرین‌ها}
\label{sec:code-exercises}

\begin{exercise}
    یک تابع پایتون با کامنت فارسی بنویسید
    و آن را با \cmd{pcc} پیچیده کنید.
    \label{exc:code-python}
\end{exercise}

\begin{exercise}
    یک الگوریتم با عنوان فارسی و متن لاتین
    بنویسید و به آن ارجاع دهید.
    \label{exc:code-algorithm}
\end{exercise}

\begin{exercise}
    یک نمونه‌ی \lr{LaTeX} با متن فارسی بنویسید
    و از \cmd{pc} استفاده کنید.
    \label{exc:code-latex}
\end{exercise}

\begin{exercise}
    یک کد \lr{C++} با کامنت فارسی بنویسید.
    \label{exc:code-cpp}
\end{exercise}

\section{خلاصه}
\label{sec:code-summary}

در این فصل، با کد و الگوریتم در \lr{MatinBook}
آشنا شدیم:

\begin{itemize}
    \item کد با محیط \env{minted} داخل
    \env{latin} نمایش داده می‌شود.

    \item برای کپشن کد، از \cmd{codecaption}
    استفاده کنید.

    \item برای کامنت فارسی در پایتون، از
    \cmd{pcc} داخل \lr{|...|} استفاده کنید.

    \item برای متن فارسی غیرکامنت در
    \lr{LaTeX} یا \lr{bash}، از \cmd{pc} استفاده
    کنید.

    \item برای کد درون‌خطی، از \cmd{inlcode}
    استفاده کنید.

    \item الگوریتم‌ها با محیط \env{algorithm}
    نمایش داده می‌شوند.

    \item کامنت فارسی در الگوریتم با
    \cmd{Comment\{|\cmd{pc}\{...\}|\}} نوشته
    می‌شود.

    \item خطاهای رایج شامل فراموش کردن
    \env{latin}، \cmd{pcc} در پایتون، و \cmd{pc}
    در \lr{LaTeX} است.
\end{itemize}

در فصل بعدی (\cref{chap:graphics})، با نمودار
و شکل در \lr{MatinBook} آشنا می‌شوید و
یاد می‌گیرید که چگونه نمودارهای حرفه‌ای
بسازید.